\section{Pilot Study: How a Benchmark Can Mislead}
\label{sec:pilot}

Before the main study, we ran a pilot benchmark (Benchmark~1) of the three single-paradigm models in a
Google Colab notebook. Its headline result was striking: the Modular network trained about
\pilotSpeedup$\times$ faster than Sequential and Parallel on every task\ifcompact\else\ (\cref{tab:pilot})\fi,
and its accuracy was only marginally worse. Taken at face value, this suggests that modularity is a free
latency win. Auditing the pilot showed that the result has nothing to do with architecture. We report the
audit because each flaw is common in small-scale benchmarking~\cite{lipton2019troubling,bouthillier2021accounting}.

\ifcompact\else  % dropped in the 6-page build
\begin{table}[t]
\centering
\caption{Pilot (Benchmark~1) results as originally reported: training time (s) and final \emph{training}
MSE, single seed. Compare the MSEs with the target variances: 34.7, 228.8, and 10.0.}
\label{tab:pilot}
\setlength{\tabcolsep}{3pt}
\begin{tabular}{@{}lrrrrrr@{}}
\toprule
 & \multicolumn{2}{c}{\textbf{Linear}} & \multicolumn{2}{c}{\textbf{Curved}} & \multicolumn{2}{c}{\textbf{Chained}}\\
\cmidrule(lr){2-3}\cmidrule(lr){4-5}\cmidrule(l){6-7}
\textbf{Model} & time & MSE & time & MSE & time & MSE\\
\midrule
\tablebody{tables/tab_pilot_body.tex}
\bottomrule
\end{tabular}
\end{table}
\fi

\textbf{Confound 1: framework, not architecture.} Sequential and Parallel were built in
Keras/TensorFlow~\cite{abadi2016tensorflow,chollet2015keras} using \texttt{model.fit}, while Modular was a
hand-written PyTorch loop. Modular and Sequential compute the same 1--128--1 function, so the
$\sim$11\,s gap is entirely framework and training-loop overhead. Framework-level differences of this size
are well documented~\cite{shi2016benchmarking}. They show that any timing comparison across frameworks
says little about architecture.

\textbf{Confound 2: uninformative inputs.} The pilot trained on $x\sim\mathcal U(0,1)$ drawn
independently of the targets, which were computed on the evenly spaced grid. With no information in the
inputs, the best any model can do is predict the target mean, which gives an MSE equal to the target variance.
The reported losses (37.8, 240, 10.4) sit just above those variances (34.7, 228.8, 10.0). In other words,
none of the pilot models learned the task, so their accuracy ranking is noise.

\textbf{Confound 3: no held-out data, one seed.} Losses were final-epoch \emph{training} losses from a single
seed with shuffling disabled, so neither generalization nor run-to-run variance could be assessed.

The main study (\cref{sec:method}) fixes all three: one framework and one training loop, true inputs, a held-out
split, and five seeds. The identical Sequential/Modular pair stays in as a built-in check. If the protocol is
sound, the two must agree on accuracy and differ only by noise.
